% !TeX root = ../../infdesc.tex

% Styles for bordered theorem environments using ntheorem and tikz packages
% Adapted from Stefan Kottwitz's answer from 2010-07-24 19:30 on the
% LaTeX.org forum: https://latex.org/forum/viewtopic.php?p=36710



% Theorems, lemmas, etc.
\tikzstyle{thmbox} = [rectangle, fill=thmbgcol, inner xsep=4pt]
\newcommand\thmbox[1]{%
  \noindent\begin{tikzpicture}%
  \node [thmbox] (box){%
    \begin{minipage}{\linewidth}%
      #1%
    \end{minipage}%
  };%
  \draw[color=thmbdcol, line width=1.5pt] (box.south west) -- (box.north west) ;
  \end{tikzpicture}}

% Definitions, notation, etc.
\tikzstyle{defbox} = [rectangle, fill=defbgcol, inner xsep=4pt]
\newcommand\defbox[1]{%
  \noindent\begin{tikzpicture}%
  \node [defbox] (box){%
    \begin{minipage}{\linewidth}%
      #1%
    \end{minipage}%
  };%
  \draw[color=defbdcol, line width=1.5pt] (box.south west) -- (box.north west) ;
  \end{tikzpicture}}

% Examples
\tikzstyle{exbox} = [rectangle, inner xsep=4pt]
\newcommand\exbox[1]{%
  \noindent\begin{tikzpicture}%
  \node [exbox] (box){%
    \begin{minipage}{\linewidth}%
      #1%
    \end{minipage}%
  };%
  \draw[color=exbdcol, line width=1.5pt] (box.south west) -- (box.north west) ;
  \end{tikzpicture}}

% Extracts
\tikzstyle{xtrbox} = [rectangle, inner xsep=4pt]
\newcommand\xtrbox[1]{%
  \noindent\begin{tikzpicture}%
  \node [xtrbox] (box){%
    \begin{minipage}{\linewidth}%
      #1%
    \end{minipage}%
  };%
  \draw[color=exbdcol, line width=1.5pt] (box.south west) -- (box.north west) ;
  \end{tikzpicture}}

% In-chapter exercises
\tikzstyle{prbox} = [rectangle, inner xsep=4pt]
\newcommand\prbox[1]{%
  \noindent\begin{tikzpicture}%
  \node [prbox] (box){%
    \begin{minipage}{\linewidth}%
      #1%
    \end{minipage}%
  };%
  \draw[color=prbdcol, line width=1.5pt] (box.south west) -- (box.north west) ;
  \end{tikzpicture}}

% Tips, remarks, etc.
\tikzstyle{tipbox} = [rectangle, fill=tipbgcol, inner xsep=4pt]
\newcommand\tipbox[1]{%
  \noindent\begin{tikzpicture}%
  \node [tipbox] (box){%
    \begin{minipage}{\linewidth}%
      #1%
    \end{minipage}%
  };%
  \draw[color=tipbdcol, line width=1.5pt] (box.south west) -- (box.north west) ;
  \end{tikzpicture}}

% solution
\tikzstyle{solbox} = [rectangle, fill=solbgcol, inner xsep=4pt]
\newcommand\solbox[1]{%
  \noindent\begin{tikzpicture}%
  \node [solbox] (box){%
    \begin{minipage}{\linewidth}%
      #1%
    \end{minipage}%
  };%
  \draw[color=solbdcol, line width=1.5pt] (box.south west) -- (box.north west) ;
  \end{tikzpicture}}

% Proofs
\tikzstyle{pfbox} = [rectangle, inner xsep=4pt]
\newcommand\pfbox[1]{%
  \noindent\begin{tikzpicture}%
  \node [pfbox] (box){%
    \begin{minipage}{\linewidth}%
      #1%
    \end{minipage}%
  };%;
  \draw[color=pfbdcol, line width=0.75pt] (box.south west) -- (box.north west) ;
  \end{tikzpicture}}

\newcommand{\setuptheoremwithborder}[4]{%
    \theoremheaderfont{\color{#1}\fontfamily{bch}\bfseries\hspace{-#3-2pt}#2\hspace{2pt}\vspace{2pt}}
    \theorembodyfont{\normalfont}
    \theorempreskip{0pt}
    \theorempostskip{0pt}
    \theorempostwork{\vspace{-10pt}}
    \theoremsymbol{}
    \let\theoremframecommand#4
}%

\newcommand{\setuptheoremwithoutborder}[3]{%
    \theoremheaderfont{\color{#1}\fontfamily{bch}\bfseries\hspace{-#3}#2\vspace{2pt}}
    \theorembodyfont{\normalfont}
    \theoremsymbol{}
    \theorempreskip{10pt}
    \theorempostskip{10pt}
}%


% Modifications to the standard 'ntheorem' package theorem styles
\makeatletter
\newtheoremstyle{infdesc}%
  {\item[\rlap{\vbox{\hbox{\hskip\labelsep \theorem@headerfont
          ##1\ ##2\theorem@separator}\hbox{\strut}}}]}%
  {\item[\rlap{\vbox{\hbox{\hskip\labelsep \theorem@headerfont
          ##1\ ##2\ $\cdot$ \itshape\mdseries ##3\theorem@separator}\hbox{\strut}}}]}
\newtheoremstyle{nonumberinfdesc}%
  {\item[\rlap{\vbox{\hbox{\hskip\labelsep \theorem@headerfont
          ##1\theorem@separator}\hbox{\strut}}}]}%
  {\item[\rlap{\vbox{\hbox{\hskip\labelsep \theorem@headerfont
          ##1\ $\cdot$ \itshape\mdseries ##3\theorem@separator}\hbox{\strut}}}]}
\newtheoremstyle{infdescpf}%
  {\item[\rlap{\vbox{\hbox{\hskip\labelsep \theorem@headerfont
          ##1\theorem@separator}\hbox{\strut}}}]}%
  {\item[\rlap{\vbox{\hbox{\hskip\labelsep \theorem@headerfont
          ##1\ \itshape\mdseries ##3\theorem@separator}\hbox{\strut}}}]}
\newtheoremstyle{nonumberinfdescpf}%
  {\item[\rlap{\vbox{\hbox{\hskip\labelsep \theorem@headerfont
          ##1\theorem@separator}\hbox{\strut}}}]}%
  {\item[\rlap{\vbox{\hbox{\hskip\labelsep \theorem@headerfont
          ##1\ ##3\theorem@separator}\hbox{\strut}}}]}
\newtheoremstyle{infdescchapex}%
  {\item[\hskip\labelsep \theorem@headerfont ##2.]}%
  {\item[\hskip\labelsep \theorem@headerfont ##2\ \itshape\mdseries(##3).]}
\makeatother

% Commands for declaring new environments.
% For example \newcommand{scholium}{scholia}{Scholium}{Scholia} creates
% a family of new theorem environments, in particular those named 'scholium'
% and 'scholia', with labels 'Scholium' and 'Scholia', respectively.

\newcommand{\newthmenv}[4]{%
    \theoremstyle{infdesc}
    \ifthenelse{\equal{#1}{theorem}}{%
        \newtheorem{#1}{#3}
    }{%
        \newtheorem{#1}[theorem]{#3}
    }%
    \newtheorem{#2}[theorem]{#4}
    \newtheorem{#1ac}[theorem]{#3\assumesAC{midgray}}
    \newtheorem{#2ac}[theorem]{#4\assumesAC{midgray}}
    \newtheorem{o#1}[theorem]{\optmark{#3}}
    \newtheorem{o#2}[theorem]{\optmark{#4}}

    \theoremstyle{nonumberinfdesc}
    \newtheorem{nn#1}{#3}
    \newtheorem{nn#2}{#4}
    \theoremstyle{infdesc}

    % Source: LaTeX Stack Exchange
    % https://tex.stackexchange.com/a/730164/145781 (Ulrike Fischer)
    \AddToHook{env/#1/begin}{\crefalias{theorem}{#1}}
    \AddToHook{env/#2/begin}{\crefalias{theorem}{#2}}
    \AddToHook{env/#1ac/begin}{\crefalias{theorem}{#1ac}}
    \AddToHook{env/#2ac/begin}{\crefalias{theorem}{#2ac}}
    \AddToHook{env/o#1/begin}{\crefalias{theorem}{o#1}}
    \AddToHook{env/o#2/begin}{\crefalias{theorem}{o#2}}

    \Crefname{#1}{#3}{#4}
    \Crefname{#2}{#3}{#4}
    \Crefname{#1ac}{#3}{#4}
    \Crefname{#2ac}{#3}{#4}
    \Crefname{o#1}{#3}{#4}
    \Crefname{o#2}{#3}{#4}
}

\newcommand{\newthmenvshaded}[4]{%
    \theoremstyle{infdesc}
    \ifthenelse{\equal{#1}{theorem}}{%
        \newshadedtheorem{#1}{#3}
    }{%
        \newshadedtheorem{#1}[theorem]{#3}
    }%
    \newshadedtheorem{#2}[theorem]{#4}
    \newshadedtheorem{#1ac}[theorem]{#3\assumesAC{midgray}}
    \newshadedtheorem{#2ac}[theorem]{#4\assumesAC{midgray}}
    \newshadedtheorem{o#1}[theorem]{\optmark{#3}}
    \newshadedtheorem{o#2}[theorem]{\optmark{#4}}

    \theoremstyle{nonumberinfdesc}
    \newshadedtheorem{nn#1}{#3}
    \newshadedtheorem{nn#2}{#4}
    \theoremstyle{infdesc}

    % Source: LaTeX Stack Exchange
    % https://tex.stackexchange.com/a/730164/145781 (Ulrike Fischer)
    \AddToHook{env/#1/begin}{\crefalias{theorem}{#1}}
    \AddToHook{env/#2/begin}{\crefalias{theorem}{#2}}
    \AddToHook{env/#1ac/begin}{\crefalias{theorem}{#1ac}}
    \AddToHook{env/#2ac/begin}{\crefalias{theorem}{#2ac}}
    \AddToHook{env/o#1/begin}{\crefalias{theorem}{o#1}}
    \AddToHook{env/o#2/begin}{\crefalias{theorem}{o#2}}

    \Crefname{#1}{#3}{#4}
    \Crefname{#2}{#3}{#4}
    \Crefname{#1ac}{#3}{#4}
    \Crefname{#2ac}{#3}{#4}
    \Crefname{o#1}{#3}{#4}
    \Crefname{o#2}{#3}{#4}
}



%
% Declarations of new theoren environments
%

%% Theorems, lemmas, etc.

% With border
\setuptheoremwithborder{thmcol}{{\color{thmbdcol}\thmintrosymbol}}{\thmsymbolwidth}{\thmbox}

\newthmenvshaded{theorem}{theorems}{Theorem}{Theorems}
\newthmenvshaded{axiom}{axioms}{Axiom}{Axioms}
\newthmenvshaded{construction}{constructions}{Construction}{Constructions}
\newthmenvshaded{conjecture}{conjectures}{Conjecture}{Conjectures}
\newthmenvshaded{corollary}{corollaries}{Corollary}{Corollaries}
\newthmenvshaded{lemma}{lemmas}{Lemma}{Lemmas}
\newthmenvshaded{proposition}{propositions}{Proposition}{Propositions}
\newthmenvshaded{assumption}{assumptions}{Assumption}{Assumptions}


%% Definitions, notation, etc.

\setuptheoremwithborder{defcol}{{\color{defbdcol}\defintrosymbol}}{\defsymbolwidth}{\defbox}

\newthmenvshaded{definition}{definitions}{Definition}{Definitions}
\newthmenvshaded{predefinition}{predefinitions}{Definition}{Definitions}
\newthmenvshaded{notation}{notations}{Notation}{Notations}
\newthmenvshaded{vocabulary}{vocabularies}{Vocabulary}{Vocabularies}



%% Examples

\setuptheoremwithborder{excol}{{\color{exbdcol}\exintrosymbol}}{\exsymbolwidth}{\exbox}

\newthmenvshaded{example}{examples}{Example}{Examples}



%% Extracts

\setuptheoremwithborder{excol}{{\color{exbdcol}\quoteintrosymbol}}{\quotesymbolwidth}{\xtrbox}

\newthmenvshaded{extract}{extracts}{Extract}{Extracts}


%% In-chapter exercises

\setuptheoremwithborder{prcol}{{\color{prbdcol}\printrosymbol}}{\prsymbolwidth}{\prbox}

\newthmenvshaded{exercise}{exercises}{Exercise}{Exercises}
\newthmenvshaded{problem}{problems}{Problem}{Problems}
\newthmenvshaded{discussion}{discussions}{Discussion}{Discussions}
\newthmenvshaded{question}{questions}{Question}{Questions}



%% Remarks, tips, etc.

% With border
\setuptheoremwithborder{tipcol}{{\color{tipbdcol}\tipintrosymbol}}{\tipsymbolwidth}{\tipbox}

\newthmenvshaded{strategy}{strategies}{Strategy}{Strategies}
\newthmenvshaded{writingprinciple}{writingprinciples}{Writing Principle}{Writing Principles}
\newthmenvshaded{remark}{remarks}{Remark}{Remarks}
\newthmenvshaded{convention}{conventions}{Convention}{Conventions}

\newshadedtheorem{deadlysin}{Deadly Sin}
\renewcommand*{\thedeadlysin}{the \Ordinalstring{deadlysin}}


\theoremstyle{nonumberinfdesc}

% solution
\setuptheoremwithborder{solcol}{{\color{solcol}\tipintrosymbol}}{\tipsymbolwidth}{\solbox}
\newshadedtheorem{mysolution}{Solution}

% proof
\setuptheoremwithborder{solcol}{{\color{solcol}\tipintrosymbol}}{\tipsymbolwidth}{\solbox}
\theoremsymbol{{\color{pfcol}$\Box$}}

\newshadedtheorem{myproof}{Proof}

\newshadedtheorem{aside}{Aside}
\newshadedtheorem{problemtip}{Problem-solving tip}
\newshadedtheorem{prooftip}{Proof tip}
\newshadedtheorem{writingtip}{Writing tip}
\newshadedtheorem{latextip}{{\LaTeX} tip}
\newshadedtheorem{formattip}{Formatting tip}
\newshadedtheorem{commonerror}{Common error}
\theoremstyle{infdesc}



%% Proofs

\setuptheoremwithborder{pfcol}{\tipintrosymbol}{\tipsymbolwidth}{\pfbox}
\theoremheaderfont{\color{pfcol}\fontfamily{bch}\bfseries\itshape}
\theoremsymbol{{\color{pfcol}$\Box$}}

\theoremstyle{nonumberinfdescpf}
\newshadedtheorem{proof}{Proof}
\newshadedtheorem{sketch}{Sketch of proof}
\newshadedtheorem{verification}{Verification}

\theoremsymbol{}
\newshadedtheorem{idea}{Idea of proof}
\theoremstyle{infdesc}



%% End-of-chapter exercises

\theoremstyle{infdescchapex}
\theoremheaderfont{\color{prcol}\fontfamily{bch}\bfseries}
\theorempreskip{5pt}
\theorempostskip{5pt}
\theorempostwork{}
\newtheorem{chapex}[theorem]{Q}
\AddToHook{env/chapex/begin}{\crefalias{theorem}{chapex}}
\Crefname{chapex}{Question}{Questions}



%% Uniformly change numbering depth

\newcommand{\numberitemswithin}[1]{\numberwithin{theorem}{#1}}
\numberitemswithin{section}



% Restating theorems without invoking the thmtools package
% Source: LaTeX Stack Exchange
% https://tex.stackexchange.com/a/51299/145781 (yo')

\newcommand{\torestate}[2]{\Copy{#1COPY}{#2}}

\newcommand{\restate}[3][]{%
\begingroup
\def\thetheorem{\ref*{#3}}
\ifthenelse{\equal{#1}{}}{\begin{#2}}{\begin{#2}[#1]}
\Paste{#3COPY}
\end{#2}%
\addtocounter{theorem}{-1}
\endgroup
}




% Hints and solutions to exercises
% Source: LaTeX Stack Exchange
% https://tex.stackexchange.com/q/415984 (Clive Newstead)
% https://tex.stackexchange.com/a/415986 (user31729)

\newcommand{\hintref}[1]{{\color{excol}\fontfamily{bch}\bfseries Hint for \Cref{#1}}\\}
\newcommand{\hintrefrange}[2]{{\color{excol}\fontfamily{bch}\bfseries Hint for \Crefrange{#1}{#2}}\\}
\newcommand{\solutionref}[1]{{\color{excol}\fontfamily{bch}\bfseries Solution to \Cref{#1}}\\}

\newwrite\hintsfile
\newwrite\solutionsfile
\AtBeginDocument{%
  % Automatically open the file at the beginning of the document
  \immediate\openout\hintsfile=\jobname.hnt
  \immediate\openout\solutionsfile=\jobname.sol
}

\NewEnviron{backhint}{%
  \immediate\write\hintsfile{%
    \expandafter\unexpanded\expandafter{\BODY}^^J%
  }%
}

\NewEnviron{backsolution}{%
  \immediate\write\solutionsfile{%
    \expandafter\unexpanded\expandafter{\BODY}^^J%
  }%
}

\newcommand{\printhints}{%
  % Closing the file
  \immediate\closeout\hintsfile% 
  \InputIfFileExists{\jobname.hnt}{}{}
}

\newcommand{\printsolutions}{%
  % Closing the file
  \immediate\closeout\solutionsfile% 
%  \InputIfFileExists{\jobname-solutions.cll}{}{}
  \InputIfFileExists{\jobname.sol}{}{}
}

\newcommand{\hint}[2]{
\label{hintorigin:#1}%
\begin{backhint}
\label{hint:#1}%
\hintref{hintorigin:#1}%
#2%
\end{backhint}%
}

\newcommand{\hintrange}[3]{
\begin{backhint}
\label{hint:#1:#2}%
\hintrefrange{#1}{#2}%
#3%
\end{backhint}%
}

\newcommand{\solution}[2]{
\label{solutionorigin:#1}
\begin{backsolution}
\label{solution:#1}
\solutionref{solutionorigin:#1}%
#2%
\end{backsolution}%
}

\newcommand{\hintsection}[1]{%
\begin{backhint}
\subsection*{#1}
\end{backhint}%
}

\newcommand{\solutionsection}[1]{%
\begin{backsolution}
\subsection*{#1}
\end{backsolution}%
}